\scanpage{21}
\deleted{
\deleted{\textcircled{\unsure{3}}}\textcircled{5}
Soient $X$ un espace topologique, $I$ un ens.
ordonné \deleted{\unsure{E}} fini,
\[
(15)\qquad (X_i)_{i\in I}
\]
une famille \inserted{croissante} de parties fermées de $X$.
\deleted{On suppose} (stratification de type $I$).
\deleted{On suppose
\textcircled{a} La famille des $X_i$ est loc$^t$ finie
\textcircled{b} $i\leq j\longrightarrow X_i\subset X_j$}
\deleted{que}
\deleted{On posera, pour $\forall i\in I$
\[
(16)\qquad X_i^*=X_i
\]}
Pour faire \deleted{des} intervenir
dans nos constructions,
on suppose donné de plus un sous-espace
bord \deleted{\unsure{de}} $\dot X$ de $X$, \uline{strictement} compatible
avec la famille des $X_i$. On veut
définir la notion de \uline{cônification}
\inserted{(resp. avec $\dot X$)} de
la stratification). Si $I$ est vide, c'est une
structure vide. Si $I=\{i\}$, $X_i=Y$, c'est
par définition un \deleted{\unsure{un}} voisinage
\deleted{tubulaire}\inserted{conique}
\inserted{strict}
de $Y$ dans $X$, \deleted{\unsure{compatible}}
tel que \inserted{la partie-bord} $\dot X$ soit \deleted{strict}
compatible avec celui-ci, que
$\dot X\cap\dot T\inserted{\overset{\mathrm{déf}}=\ddot T}$
soit un sous-espace bord de $\dot T$, qu'on ait
$\deleted{0.(\dot X\cap\dot T)}
0.\ddot T=\underbrace{\dot X\cap Y}_{\dot Y}$
(\unsure{a priori on} \deleted{\unsure{équitable}}\inserted{soit}),
\deleted{\unsure{que on}} $\ddot T$ soit un
sous-espace bord de $Y$, et que $\dot T\longrightarrow Y$ soit
dans un sens à préciser.
}
\margin{20}
